\section{Evidence Synthesis by Quality Dimension}
\label{sec:dimensions}

This section will contain the evidence-driven synthesis of the academic corpus after the peer-review audit and final classification freeze. Each subsection will follow the same structure: definition and scope, dominant quality constructs, evaluation methods, representative studies, recurring datasets or benchmarks, evidence of cross-stage relations, limitations, and open gaps.

\subsection{Input and Specification Quality}
\label{sec:input-quality}

% TODO:
% - Prompt sensitivity and robustness.
% - Instruction/specification quality.
% - Prompt optimization and perturbation.
% - Distinguish prompt use from prompt-quality evaluation.
% - Synthesize evidence linking input variation to downstream output quality.

\subsection{Context and Retrieval Quality}
\label{sec:context-quality}

% TODO:
% - Retrieval relevance, recall, precision, grounding, context sufficiency.
% - RAG evaluation frameworks.
% - Separate retrieval-quality metrics from final-answer metrics.
% - Analyze retrieval -> output relationships and cases evaluated only black-box.

\subsection{Output and Response Quality}
\label{sec:output-quality}

% TODO:
% - Correctness, factuality, faithfulness, relevance, completeness, helpfulness.
% - Code and task-specific output quality.
% - Reference-based vs reference-free evaluation.
% - Human, learned, semantic, and judge-based evaluation.
% - Safety/fairness/robustness as cross-cutting output properties.

\subsection{Conversation and Interaction Quality}
\label{sec:conversation-quality}

% TODO:
% - Multi-turn consistency and state maintenance.
% - User-goal completion and adaptation to feedback.
% - Dynamic evaluation and user simulation.
% - Long-horizon failure propagation.

\subsection{Agent, Process, and Trajectory Quality}
\label{sec:agent-quality}

% TODO:
% - Planning, decomposition, routing, tool calls, trajectory quality.
% - Stepwise vs final-response evaluation.
% - Trace-aware evaluation.
% - Multi-agent coordination and protocol-level behavior.
% - Task success vs process correctness.

\subsection{Operational and System Quality}
\label{sec:operational-quality}

% TODO:
% - Latency, throughput, cost/query, tokens, memory, energy.
% - Availability, observability, privacy, safety.
% - Development vs production evaluation.
% - Quality-cost trade-offs.
% - Relationship between semantic quality and deployment constraints.

\subsection{Evaluator Reliability and Meta-Evaluation}
\label{sec:evaluator-reliability}

% TODO:
% - Human agreement and evaluator alignment.
% - LLM-judge consistency, bias, calibration, prompt sensitivity.
% - Multi-judge dependence and ensemble reliability.
% - Distinguish use of a judge from validation of a judge.
% - Statistical reliability measures and meta-evaluation studies.
